\documentclass{article}
\usepackage[utf8]{inputenc}
\usepackage[ngerman]{babel} % deutsch ahh
\usepackage{amsmath}
\usepackage{amsfonts}
\usepackage[colorlinks=true, urlcolor=blue, linkcolor=blue]{hyperref}

\title{Mögliche Lösungen zum Reserve Mathe 5 BAC 2017 A-Teil}
\author{Milix244}
\date{\today}

\begin{document}

\maketitle

\section*{Aufgabe 1}
Berechnen Sie
\[\int(x-2)e^x\ dx\]

\subsection*{Lösung}
Ansatz ist die partielle Integration \\
Formel:
\[\int u'\cdot v\ dx=u\cdot v-\int u\cdot v'\ dx\]
Wahl von u' und v (LIATE):
\[u'(x)=e^x\]
\[v(x)=(x-2)\]
Einsetzen:
\[\int e^x\cdot (x-2)\ dx\]
Wobei:
\[u(x)=e^x\]
\[v'(x)=1\]
Daraus folgt:
\[=(e^x)\cdot (x-2)-\int e^x\cdot 1\ dx\]
\[=(e^x)\cdot (x-2)-(e^x)\]
\[=(e^x)\cdot ((x-2)-1)\]
\[=(e^x)\cdot (x-3)\]
C hinzufügen:
\[=(e^x)\cdot (x-3)+C,\ (C\in\mathbb{R})\]

\newpage
\section*{Aufgabe 2}
Gegeben ist die Ebene \(\pi :2x-2y-z=0\). \\
Ermitteln Sie eine Koordinatengleichung für jede Ebene, die parallel zur
Ebene \(\pi\) ist und von \(\pi\) den Abstand 5 besitzt.
\subsection*{Lösung}
Gegeben ist die Koordinatenform der Ebene. \\
Gesucht sind zwei parallele Ebenen (die durch eine Gleichung beschrieben wird \(E_{parallel}\)) von Ebene \(\pi\) mit Abstand \(d=5\). \\
Gegebene Formel (Punkt - Ebene, Zur Übersicht wird \(D\left(M,\pi \right)\) anstatt von \(d\left(M,\pi \right)\) verwendet):
\[D\left(M,\pi \right)=\frac{\left|\overrightarrow{AM}\cdot \vec{n}\right|}{\left|\vec{n}\right|}=\frac{\left|ax_{M}+by_{M}+cz_{M}+d\right|}{\sqrt{a^{2}+b^{2}+c^{2}}}\]
Umformung von \(\pi\):
\[\pi :2x-2y-z=0\]
Wobei:
\[\vec{n}=\begin{pmatrix} 2 \\ -2 \\ -1 \end{pmatrix}\]
Also:
\[\pi :\begin{pmatrix} 2 \\ -2 \\ -1 \end{pmatrix}\cdot \left(\overrightarrow{X}-\overrightarrow{P_0}\right)=0\]
\[\pi :\begin{pmatrix} 2 \\ -2 \\ -1 \end{pmatrix}\cdot\overrightarrow{X}-\begin{pmatrix} 2 \\ -2 \\ -1 \end{pmatrix}\cdot\overrightarrow{P_0}=0\]
Wobei:
\[\vec{n}\cdot\overrightarrow{P_0}=d\]
und:
\[d=0\]
Erhält man:
\[\pi :\begin{pmatrix} 2 \\ -2 \\ -1 \end{pmatrix}\cdot \overrightarrow{X}=0\]
Da man in diesen Fall parallele Ebenen sucht, ist der Normalenvektor \(\vec{n}\) bei den gesuchten Ebenen wie bei Ebene \(\pi\) gleich. \\
Die gesuchten Ebenen haben die Form:
\[E_{\text{parallel}}: 2x-2y-z=d'\]
wobei \(d'\) so bestimmt werden muss, dass der Abstand zwischen \(\pi\) und \(E_{\text{parallel}}\) gleich 5 ist (\(D\left(M,\pi \right)=5\)). \\
Man leitet sich folgende Formel her (M ist ein Punkt auf der \(E_{parallel}\)):
\[D\left(M,\pi \right)=\frac{\left|\overrightarrow{AM}\cdot \vec{n}\right|}{\left|\vec{n}\right|}\]
\[=\frac{\left|\left(\overrightarrow{M}-\overrightarrow{A}\right)\cdot \vec{n}\right|}{\left|\vec{n}\right|}\]
\[=\frac{\left|\overrightarrow{M}\cdot \vec{n}-\overrightarrow{A}\cdot \vec{n}\right|}{\left|\vec{n}\right|}\]
Wobei:
\[\overrightarrow{M}\cdot \vec{n}=d'\]
\[\overrightarrow{A}\cdot \vec{n}=d\]
Demnach:
\[D\left(M,\pi \right)=\frac{\left|d'-d\right|}{\left|\vec{n}\right|}\]
\[D\left(M,\pi \right)=\frac{\left|d'-d\right|}{\sqrt{x^2+y^2+z^2}}\]
Mit der Formel:
\[5=\frac{|d'-d|}{\sqrt{2^2+(-2)^2+(-1)^2}}\]
\[5=\frac{|d'-0|}{\sqrt{4+4+1}}=\frac{|d'|}{3}\]
\[5\cdot 3=|d'|\]
\[15=|d'|\]
Also \(d'=15\) oder \(d'=-15\). \\
Die beiden parallelen Ebenen sind:
\[E_{parallel}: 2x-2y-z=\pm 15\]
Also:
\[E_{1}: 2x-2y-z=15,\quad E_{2}: 2x-2y-z=-15\]

\newpage
\section*{Aufgabe 3}
Ein idealer Würfel mit den Augenzahlen 1 bis 6 wird 5 Mal geworfen. \\
Berechnen Sie die Wahrscheinlichkeit, dass die Augenzahl mindestens
4 Mal eine gerade Zahl ist.
\subsection*{Lösung}

\newpage
\section*{Aufgabe 4}
Die Folgen (\(u_n\)) und (\(v_n\)) sind gegeben durch
\[\begin{cases}
        u_1 = 5 \quad \text{und} \quad v_1 = 2 \\[2mm]
        u_{n+1} = u_n + v_n \quad \text{und} \quad v_{n+1} = \dfrac{u_{n+1} + 2v_n}{2}, \quad n \geq 1.
    \end{cases}\]
Nun betrachten wir die Folge (\(w_n\)) mit \(w_n=u_n-v_n\). \\
Zeigen Sie, dass (\(w_n\)) eine geometrische Folge ist.
\subsection*{Lösung}

\newpage
\section*{Aufgabe 5}
Die komplexe Zahl \(z=-2\sqrt{3}+2i\) ist eine der vierten Wurzeln einer
komplexen Zahlen \(w\). \\
Bestimmen Sie die anderen vierten Wurzeln von \(w\) und stellen Sie alle
vierten Wurzeln von \(w\) in der Gaußschen Zahlenebene dar.
\subsection*{Lösung}

\newpage
\section*{Aufgabe 6}
Von einer Funktion \(f\) ist folgendes bekannt:
\[f(4)=3\quad\text{und}\quad\int^4_0f(t)dt=13\]
(Deutsche BAC Version hat ein Übersetzungsfehler mit \(f(4)=4\) in der Aufgabenstellung) \\
Eine Funktion \(F\) ist gegeben durch \(F(x)=\int_{0}^{x}f(t)dt\). \\
Bestimmen Sie eine Gleichung für die Tangente am Schaubild von \(F\)
im Punkt mit der Abszisse \(x=4\).
\subsection*{Lösung}

Gesucht ist die Tangentengleichung an \(x=4\) von \(F(x)\). Wir wissen, dass \(f(x)\) die Ableitung von \(F(x)\) ist. \\
Daraus folgt bereits:
\[\text{Tangentengleichung: }y=mx+t\]
\[m=3\]
\[y=3\cdot x+t\]
Zusätzlich brauchen wir ein Punkt auf \(F(x)\). Gegeben ist:
\[F(4)=\int^4_0f(t)dt=13\]
\[P(4|13)\]
Einsetzen:
\[13=3\cdot 4+t\]
\[13=12+t\]
\[1=t\]
Daraus folgt die Tangentengleichung:
\[y=3x+1\]
(Die deutsche Übersetzung wäre entsprechend: \(y=4x-3\))

\newpage
\section*{Aufgabe 7}
Ein Sack enthält 4 weiße und \(n\) rote Kugeln. \\
Zwei Kugeln werden nacheinander ohne Zurücklegen entnommen. \\
Die Wahrscheinlichkeit dass beide entnommenen Kugeln rot sind,
beträgt \(\frac{1}{3}\). \\
Bestimmen Sie den Wert von \(n\).
\subsection*{Lösung}

\newpage
\begin{center}
    {\LARGE \textbf{Credits}}
\end{center}
\vspace{1cm}
\noindent\textbf{Main Editor: } Milix
\vspace{0.5cm}
\begin{center}
    \href{https://github.com/Milix244}{GITHUB PROFILE}  \\
    \vspace{0.3cm}
    \href{https://github.com/Milix244/Reserve-BAC-2017}{GITHUB REPO}  \\
    \vspace{0.3cm}
    \href{https://discord.gg/EpQ3eukxJ4}{DISCORD SERVER}
\end{center}
\vspace{1cm}
\textit{\copyright\ 2026 -- Alle Rechte vorbehalten trust.}


\end{document}